%% TF de l'impulsion decalee delta(t-t0) : spectre complexe de module
%% constant 1, partie reelle cos, partie imaginaire -sin. Repris du
%% poly (section 3). Consomme par slides.
\begin{tikzpicture}[x=1.2cm,y=1cm]
    \draw[-Latex] (-2,0) -- (2,0) node[right] {$t$};
    \draw[niceblue,ultra thick,-Latex] (1,0) -- (1,1.6) node[above,yshift=2pt] {$\delta(t-t_0)$};
    \node[below] at (1,0) {$t_0$};
    \node[below left] at (0,0) {$0$};
\end{tikzpicture}
\hspace{1cm}
\begin{tikzpicture}
    \begin{axis}[
        width=6.5cm, height=5cm,
        xlabel={$f$},
        ylabel={\color{green!60!black} $Re(X(f))$,\color{red} $Im(X(f))$ },
        xmin=-3, xmax=3, ymin=-1.3, ymax=1.9,
        axis lines=middle,
        xlabel style={at=(current axis.right of origin), anchor=west},
    ]
    \addplot[domain=-3:3, samples=200, green!60!black, thick] {cos(deg(2*3.14159*x))};
    \addplot[domain=-3:3, samples=200, red, thick] {-sin(deg(2*3.14159*x))};
    \addplot[domain=-3:3, samples=2, niceblue, dashed, thick] {1};
    % \addplot[domain=-3:3, samples=2, niceblue, dashed, thick] {-1};
    \end{axis}
\end{tikzpicture}
